% Desarrollo acotado: lector energético intervalar del registro K.
% Estatuto: formalización añadida sobre K y sobre las identidades de Schur
% ya presentes en el corpus. No se atribuye novedad histórica a Schur.
\section{Formes de Dirichlet et équivalence variationnelle du raffinement}
\label{sec:energia-refinamiento-k}

La représentation positive du registre dodécaphasique définit une partition ordonnée au moyen de
\[
 K=(234,543,140,729,659,824,621,58,914,794,146,601),
 \qquad d_j=K_j/6263,
 \qquad t_i=\sum_{j=1}^id_j.
\]
Le registre est reçu de la composition
\(x_{\rm term}\mapsto\mathcal R_{12}\mapsto
(b^{90},b^{120},Q_{\rm TPK})\mapsto U_{\rm sgn}\mapsto K\),
postérieure à APP--TRIT--TPK et à son état enrichi, conformément aux lecteurs
reconstruits dans \cite{hmtX}. Cette section applique
un lecteur quadratique à la partition résultante. La comparaison entre
profondeurs s’effectuera par restrictions aux extrémités héritées, de sorte
que chaque coefficient de l’énergie conserve son origine dans une séparation
positive déjà produite.

Le corpus contient une énergie pondérée de différences pour le graphe des
deux mémoires et une réduction de Schur des formes complètes. On
réunit ici ces opérations sur le réseau d’intervalles du registre. La
spécialisation aux conductances \(1/d_j\), l’opérateur d’interpolation
et sa compatibilité avec la déformation des séparations constituent
la formalisation développée dans cette section.

\subsection{Reconstruction des séparations à partir de la matrice de rigidité}

Soit \(d=(d_1,\ldots,d_m)\), avec \(d_j>0\), et soient
\(f=(f_0,\ldots,f_m)\in\mathbb C^{m+1}\) les valeurs d’un observable
aux extrémités de la partition. Nous définissons
\begin{equation}
 \mathcal E_d(f)=\sum_{j=1}^{m}\frac{|f_j-f_{j-1}|^2}{d_j},
 \qquad L_d=D^*\operatorname{diag}(d_1^{-1},\ldots,d_m^{-1})D,
 \label{eq:energia-intervalar}
\end{equation}
où \(D:\mathbb C^{m+1}\to\mathbb C^m\) est l'opérateur d'incidence orientée, \((Df)_j=f_j-f_{j-1}\), et l'astérisque désigne l'adjonction hermitienne. Ainsi \(\mathcal E_d(f)=f^*L_df\). La matrice de rigidité a pour entrées
\[
 (L_d)_{00}=d_1^{-1},\quad (L_d)_{mm}=d_m^{-1},\quad
 (L_d)_{ii}=d_i^{-1}+d_{i+1}^{-1}\quad(1\le i<m),
\]
\[
 (L_d)_{i-1,i}=(L_d)_{i,i-1}=-d_i^{-1},
\]
et les autres sont nuls. Elle est hermitienne positive semi-définie, son noyau
est constitué des fonctions constantes et sa forme est définie positive
sur le quotient par ce noyau. La matrice complète retrouve les
séparations par
\begin{equation}
 d_j=-\frac1{(L_d)_{j-1,j}}.
 \label{eq:recuperacion-d-rigidez}
\end{equation}

\begin{proposition}[Équivalence entre séparations et lecteur nodal complet]
\label{prop:energia-bidireccional}
L'application \(d\mapsto L_d\) est une bijection entre \((0,\infty)^m\) et l'ensemble des matrices réelles symétriques tridiagonales dont les entrées immédiatement adjacentes à la diagonale sont négatives et dont la somme de chaque ligne est nulle. Son inverse est \eqref{eq:recuperacion-d-rigidez}. La normalisation \(\sum_jd_j=1\) équivaut à \(\sum_{j=1}^m-1/(L_d)_{j-1,j}=1\).
\end{proposition}
\begin{proof}
Les formules explicites de \(L_d\) vérifient toutes les conditions.
Réciproquement, les entrées adjacentes négatives déterminent des séparations
positives uniques par \eqref{eq:recuperacion-d-rigidez}. La somme nulle de
chaque ligne détermine ses entrées diagonales comme la somme des
conductances incidentes. La tridiagonalité fixe à zéro les entrées
restantes. On récupère exactement la matrice de
\eqref{eq:energia-intervalar} ; sa semi-définie positivité découle alors
de la factorisation \(D^*\operatorname{diag}(1/d_j)D\).
\end{proof}

La dérivée discrète pondérée \(J_j=(f_j-f_{j-1})/d_j\) est le flux orienté de l'arête \(j\). Aux sommets intérieurs, \((L_df)_i=J_i-J_{i+1}\) ; aux extrémités, \((L_df)_0=-J_1\) et \((L_df)_m=J_m\). L'équation intérieure \(L_df=0\) exprime que le flux a la même valeur sur chaque arête. Cette définition est celle d'un observable quadratique du lecteur d'intervalles ; sa réalisation comme grandeur physique requiert la coordinatisation dimensionnelle pertinente.

\subsection{Interpolation harmonique, complément de Schur et reconstruction du détail}

Raffinons chaque séparation sous la forme \(d_j=\sum_{c=1}^{r_j}\delta_{j,c}\), avec \(\delta_{j,c}>0\). Soit \(\widetilde d\) la concaténation de ces séparations, \(R\) la restriction aux anciennes extrémités et \(H\) l'interpolation définie par
\begin{equation}
 (Hf)_{j,c}=(1-\theta_{j,c})f_{j-1}+\theta_{j,c}f_j,
 \qquad
 \theta_{j,c}=\frac{\sum_{a=1}^c\delta_{j,a}}{d_j}.
 \label{eq:interpolacion-armonica-k}
\end{equation}
Les extrémités partagées ne sont comptées qu’une fois. On a \(RH=I\).

\begin{theorem}[Raffinement isométrique de l’énergie]
\label{thm:energia-schur-refinamiento}
Pour tout vecteur de valeurs anciennes \(f\),
\begin{equation}
 H^*L_{\widetilde d}H=L_d,
 \qquad L_{\widetilde d}H=R^*L_d,
 \qquad
 \min_{Rg=f}\mathcal E_{\widetilde d}(g)=\mathcal E_d(f).
 \label{eq:isometria-energia-k}
\end{equation}
Le minimiseur est unique et vaut \(Hf\). Tout \(g\) admet la
décomposition unique
\begin{equation}
 g=H(Rg)+z,\qquad Rz=0,
 \qquad
 \mathcal E_{\widetilde d}(g)
 =\mathcal E_d(Rg)+\mathcal E_{\widetilde d}(z).
 \label{eq:detalle-energia-k}
\end{equation}
\end{theorem}
\begin{proof}
Sur une ancienne arête, écrivons \(h_c=g_{j,c}-g_{j,c-1}\), de sorte que \(\sum_ch_c=f_j-f_{j-1}=a\). Si \(h_c=\delta_{j,c}a/d_j+r_c\), alors \(\sum_cr_c=0\) et
\[
 \sum_c\frac{|h_c|^2}{\delta_{j,c}}
 =\frac{|a|^2}{d_j}+\sum_c\frac{|r_c|^2}{\delta_{j,c}}.
\]
L’égalité s’obtient en développant ; le terme croisé est proportionnel à
\(\sum_cr_c\) et s’annule. Le minimum est atteint exactement lorsque
tous les \(r_c\) sont nuls, condition qui détermine
\eqref{eq:interpolacion-armonica-k}. En sommant les arêtes, on obtient
le minimum et la décomposition énergétique. Le flux de \(Hf\) est
constant à l’intérieur de chaque ancienne arête, de sorte que ses composantes
intérieures de \(L_{\widetilde d}Hf\) sont nulles et ses composantes héritées
coïncident avec \(L_df\). Cela prouve la seconde identité matricielle ; en
multipliant par \(H^*\), en utilisant \(RH=I\), on obtient la première.
\end{proof}

Le terme \(\mathcal E_{\widetilde d}(z)\) est non négatif et s’annule
exactement pour \(z=0\) : un détail d’énergie nulle serait constant et,
en s’annulant aux extrémités héritées, il s’annulerait sur tout le réseau.
La paire \((Rg,z)\) conserve le vecteur raffiné complet et permet
de le retrouver par \(g=H(Rg)+z\). La minimisation sélectionne sa
composante harmonique ; l’archivage du détail \(z\) maintient
l’information que la seule forme réduite ne distingue pas. La conservation
de l’énergie observée et la conservation de l’état raffiné sont
ainsi exprimées par des applications différentes et explicites.

En ordonnant les sommets comme extrémités héritées \(E\) et sommets intérieurs \(I\), nous écrivons
\[
 L_{\widetilde d}=\begin{pmatrix}A&B^*\\B&C\end{pmatrix}.
\]
Pour un raffinement comportant des sommets intérieurs, le bloc \(C\) est défini positif : une fonction nulle aux extrémités dont l'énergie est nulle est constante sur chaque sous-intervalle et, par conséquent, identiquement nulle. On obtient
\begin{equation}
 H=\begin{pmatrix}I\\-C^{-1}B\end{pmatrix},\qquad
 \operatorname{Schur}_{I}(L_{\widetilde d})
 =A-B^*C^{-1}B=L_d.
 \label{eq:schur-intervalar-k}
\end{equation}
Si le raffinement conserve tous les sommets et n'en introduit aucun, la même affirmation se réduit à \(H=I\) et \(L_{\widetilde d}=L_d\).

\begin{proposition}[Naturalité des prolongements et de l'élimination]
\label{prop:energia-naturalidad-sucesiva}
Pour trois partitions emboîtées, les interpolations harmoniques et les compléments de Schur satisfont
\begin{equation}
 H_{n+2\leftarrow n}
 =H_{n+2\leftarrow n+1}H_{n+1\leftarrow n},
 \qquad
 \operatorname{Schur}_{n+1\to n}
 \bigl(\operatorname{Schur}_{n+2\to n+1}(L_{n+2})\bigr)=L_n.
 \label{eq:energia-naturalidad-sucesiva}
\end{equation}
\end{proposition}
\begin{proof}
L’interpolation de valeurs affines sur un intervalle demeure la
même fonction affine après toute subdivision de cet intervalle.
Cela prouve la première identité. Minimiser d’abord sur les sommets
incorporés au dernier niveau puis sur ceux du niveau intermédiaire
revient à minimiser sur leur totalité, avec les mêmes extrémités héritées fixées.
L’identité variationnelle de \eqref{eq:isometria-energia-k} et l’unicité
du minimiseur prouvent la seconde identité. La restriction du minimum
se réalise donc au moyen de la même composition que le prolongement
harmonique des données.
\end{proof}

\subsection{Réponse de Dirichlet à Neumann et défaut énergétique d’une extension}

Si seules les extrémités initiale et finale sont conservées, soit \(D_0=\sum_{j=1}^md_j\). Pour des valeurs \((a,b)\), la solution harmonique est
\[
 f_i=a+(b-a)\frac{\sum_{j=1}^id_j}{D_0},\qquad
 J_j=(b-a)/D_0.
\]
Son opérateur de Dirichlet à Neumann, qui transforme les valeurs de frontière en
flux sortants, est
\begin{equation}
 \Lambda_d=\frac1{D_0}
 \begin{pmatrix}1&-1\\-1&1\end{pmatrix},\qquad
 \min\mathcal E_d=\frac{|b-a|^2}{D_0}.
 \label{eq:dtn-intervalar-k}
\end{equation}
Lorsque la valeur est prescrite à tous les sommets hérités, leur opérateur de réponse multiponctuelle est, par le même argument, \(\Lambda=L_d\).

Une extension approchée \(g=(f,y)\), où \(y\) réunit les valeurs
aux sommets intérieurs, a pour résidu intérieur
\(r=Bf+Cy\). L’identité exacte
\begin{equation}
 \mathcal E_{\widetilde d}(g)-f^*\Lambda f
 =r^*C^{-1}r\ge0
 \label{eq:residual-dtn-k}
\end{equation}
se déduit de \(y+C^{-1}Bf=C^{-1}r\) et de la décomposition
\[
 g^*L_{\widetilde d}g
 =f^*(A-B^*C^{-1}B)f
 +(y+C^{-1}Bf)^*C(y+C^{-1}Bf).
\]
Ainsi, la valeur effective et le détail maintiennent une décomposition
hermitienne complète. En particulier, si une
extension linéaire est \(\widetilde H=\begin{pmatrix}I\\Z\end{pmatrix}\), sa matrice de résidus est \(\mathscr R=B+CZ\) et
\[
 \widetilde H^*L_{\widetilde d}\widetilde H-\Lambda
 =\mathscr R^*C^{-1}\mathscr R\succeq0.
\]
Le résidu nul caractérise l'interpolation harmonique exacte. Ce résidu énergétique se situe dans l'espace des sommets intérieurs ; l'annulation des résidus des formes méromorphes appartient à une autre application de frontière et conserve sa propre définition.

\subsection{Commutation du raffinement et du flux des séparations}

Soit \(u=P_3P_{11}K\) la direction algébrique déjà retrouvée à partir des
lecteurs du registre \cite{hmtX}. Pour le paramètre réel \(s\), définissons
\[
 d_j(s)=\frac{K_je^{su_j}}{\sum_\ell K_\ell e^{su_\ell}}.
\]
Étant fixées des fractions positives \(\vartheta_{j,c}\) de somme un pour chaque
\(j\), le raffinement hérité est
\[
 \delta_{j,c}(s)=\vartheta_{j,c}d_j(s),\qquad
 u_{j,c}=u_j.
\]
Produire d'abord les poids raffinés \(K_j\vartheta_{j,c}\) et les déformer donne la même collection : les successeurs ont une exponentielle identique et leurs coefficients ont pour somme \(K_j\). Pour tout \(s\),
\begin{equation}
 \operatorname{Schur}_{I}(L_{\widetilde d(s)})=L_{d(s)},
 \qquad H^*L_{\widetilde d(s)}H=L_{d(s)}.
 \label{eq:schur-flujo-k}
\end{equation}
Dans cette comparaison parent--enfants, \(H\) est constant : ses coefficients
sont les sommes partielles de \(\vartheta_{j,c}\). En revanche,
l’interpolation à partir des deux extrémités globales utilise
\(t_i(s)=\sum_{j\le i}d_j(s)\) et change avec le paramètre. Les deux
affirmations correspondent à des restrictions différentes du même réseau.

La normalisation \(\sum_jd_j(s)=1\) fait que l'opérateur global à deux extrémités reste égal à \(\left(\begin{smallmatrix}1&-1\\-1&1\end{smallmatrix}\right)\). La matrice de tous les nœuds \(L_{d(s)}\) conserve, elle, la déformation et la récupère au moyen de \eqref{eq:recuperacion-d-rigidez}. La perte d'information dans le lecteur à deux extrémités est identifiée exactement : il retient la longueur totale et ses flux, tandis que le lecteur nodal retient toutes les séparations. L'archive des détails fournit le complément nécessaire pour reconstruire l'observable raffiné.
